\documentclass[11pt]{article}

\usepackage[margin=1in]{geometry}
\usepackage{hyperref}
\usepackage{enumitem}
\usepackage{parskip}

\title{\textbf{Belief-Aware Reinforcement Learning in Cooperative Games and Financial Markets}}
\author{EECS 6892 Reinforcement Learning --- Project Proposal}
\date{}

\begin{document}

\maketitle

\section*{Project Description}

Reinforcement learning agents trained through self-play often perform well when interacting with identical copies of themselves but fail when paired with agents using different strategies. This problem appears in many multi-agent environments where the behavior of other participants is diverse and only partially observable. 

This project studies whether a reinforcement learning agent that maintains a \textit{belief model} about the type of agents it interacts with can outperform a standard RL agent that does not explicitly model others. The belief-aware agent observes the actions of its counterpart and produces a probability distribution over possible agent types. This belief representation is provided as additional input to the policy network so that the agent can adapt its behavior in real time.

We test this idea in two complementary environments. First, we evaluate belief-aware RL in a simulated multi-agent financial market with heterogeneous traders (noise, momentum, and mean-reversion strategies). Second, we test the same architecture in the cooperative card game Hanabi, which is widely used as a benchmark for multi-agent coordination. Using both environments allows us to evaluate whether belief modeling improves performance in both competitive and cooperative settings.

\section*{Research Directions and Expected Outcomes}

The project focuses on the following questions:

\begin{itemize}[leftmargin=*]
\item Does incorporating a belief model about other agents improve policy performance in partially observable multi-agent environments?
\item Can a shared belief architecture generalize across different domains such as financial markets and cooperative games?
\item How quickly can an agent infer the type of its counterpart from early interactions?
\end{itemize}

We will compare two agents: a baseline RL agent trained without opponent modeling and a belief-aware RL agent that jointly learns a belief module and policy network. Experiments will measure performance using environment-specific metrics such as trading profit (PnL) and Sharpe ratio in the market simulation, and cross-play score in Hanabi. 

We expect that belief-aware agents will achieve higher robustness when interacting with diverse partners or market participants, demonstrating the value of explicit belief inference in multi-agent reinforcement learning.

\section*{Reading List}

The project is based on prior work in multi-agent reinforcement learning, opponent modeling, and cooperative game learning.

\begin{itemize}[leftmargin=*]
\item Foerster et al., \textit{Learning with Opponent-Learning Awareness}, AAMAS 2018.
\item Hu and Wellman, \textit{Nash Q-Learning for General-Sum Stochastic Games}, JMLR 2003.
\item Bard et al., \textit{The Hanabi Challenge: A New Frontier for AI Research}, AI Magazine 2020.
\item Hu et al., \textit{Other-Play for Zero-Shot Coordination}, ICML 2020.
\item Carroll et al., \textit{On the Utility of Learning about Humans for Human-AI Coordination}, NeurIPS 2019.
\item Nevmyvaka et al., \textit{Reinforcement Learning for Optimized Trade Execution}, ICML 2006.
\end{itemize}

\end{document}